\documentclass[11pt, a4paper]{article}
\usepackage[a4paper, margin=2.5cm]{geometry}
\usepackage{tabularx}
\usepackage{hyperref}

\hypersetup{
    colorlinks=true,
    linkcolor=black,
    urlcolor=black
}

\begin{document}

\begin{titlepage}
    \centering
    \vspace*{2cm}

    {\LARGE \bfseries Analyzing brain tumors using MRI scans \par}
    \vspace{0.5cm}
    {\Large Initial System Specification \par}

    \vspace{2cm}

    \begin{tabularx}{\textwidth}{@{} >{\bfseries}l X @{}}
        Course:         & Software Engineering (SE)                                            \\
        Programme:      & Informatics (INF), Semester 5, Academic Year 2026/2027               \\
        Authors:        & Jakub Toporek, Zuzanna Safaryjska, Xawier Słupik, Mateusz Pyrek,     \\
                        & Maciej Gajdziński, Oskar Bednarek, Eryk Borowiecki, Michał Głogowski \\
        Created:        & 13.11.2026                                                           \\
        Modified:       &                                                                      \\
        Version:        & {[1.0.0]}                                                            \\
        Classification: & Confidential                                                         \\
    \end{tabularx}

    \vfill
\end{titlepage}

\tableofcontents
\newpage

\section{Document history}

\noindent
\renewcommand{\arraystretch}{1.3}
\begin{tabularx}{\textwidth}{|l|l|l|l|X|}
    \hline
    \textbf{Version} & \textbf{Date} & \textbf{Author} & \textbf{Approved by} & \textbf{Additional information} \\
    \hline
    1.0.0            & 2026-11-13    & JT, ZS          &                      & Initial document draft          \\ \hline
    1.1.0            & 2026-11-15    & XS, MiG         &                      & Requirements and file formats   \\
    \hline
\end{tabularx}

\section{The introduction}

\subsection{Primary Purpose}

Magnetic Resonance Imaging (MRI) is a cornerstone of non-invasive oncology, providing unmatched soft-tissue contrast for detecting, staging, and monitoring solid tumors. We will use 4D images for 3D scans with time as a fourth dimension. We distinguish 4 different types of images for instance FLAIR, T1w, T1gd, T2w.

Despite its clinical utility, 4D MRI evaluation is hampered by severe operational bottlenecks in legacy medical workflows. Diagnostic workstations predominantly rely on heavyweight, monolithic desktop software (e.g., OsiriX, Horos, 3D Slicer) installed on dedicated, high-spec clinical workstations within hospital intranets. This paradigm introduces significant friction.

\subsection{Implementation}

We decided to build a desktop app. We will be using Python with Napari framework to build robust and intuitive tool for 4D MRI visualization.

\section{Requirements}

\subsection{Functional requirements}

\begin{tabularx}{\textwidth}{|l|X|X|}
\hline
ID  & Requirement                                                   & Acceptance Criteria \\ \hline
FR1 & User can load images.                                         & Images are viewable in the browser. \\ \hline
FR2 & Users can browse through the entire scan slicing and timing.  & All imaging is visible. \\ \hline
FR3 & Users can highlight selected areas.                           & User-defined highlights appear over the images. Uses designated mouse pointer tools to highlight them. \\ \hline
FR4 & Users can get the volume of selected areas.                   & Returns to the user the volume of the selection. \\ \hline
FR5 & Users can see scan details (image dimensions, size, etc.).    & Can be inspected for any scanning via designated panel. \\ \hline
\end{tabularx}

\subsection{Non-functional requirements}

\begin{tabularx}{\textwidth}{|l|X|X|}
\hline
ID  & Requirement                       & Acceptance Criteria \\ \hline
NFR1 & Slice Scrubbing Responsiveness.  & Scrolling through 2D slices in Multi-Planar Reconstruction (MPR) viewports must maintain $\geq$ 60 FPS (sub-16 ms frame time). \\ \hline
NFR2 & Lossless Pixel Fidelity.         & The application must never apply lossy compression or unintended bit-depth truncation to raw 12-bit/16-bit MRI intensities. \\ \hline
NFR3 & Crash recoverability.            & Unsaved segmentation work must be protected against sudden process termination or power loss. In the event of an unexpected application crash, the Recovery Point Objective (RPO) is $<$ 60 seconds of annotation work. \\ \hline
NFR4 & Keyboard-driven workflow.        & Common inspection and segmentation actions must be accessible without looking through several menus. \\ \hline
NFR5 & Dataset speed.                   & Opening standard medical volumes must not bottleneck the user's workflow. Reconstructing of a file should be completed in $<$ 5s \\ \hline
\end{tabularx}

\section{File formats}
\subsection{DICOM (.dcm)}
\begin{itemize}
    \item 1 file is equivalent to 1 image slice
    \item File is a binary made up of 3 sections:
    \begin{itemize}
        \item Preamble (128 bytes + 4 bytes prefix) - first 128 bytes are for compatibility headers followed by text “DICM” (prefix).
        \item Header - collection of tags (key-value pairs) with key defined as hex value. Stores data such as Rows / Columns, Pixel Spacing, Patient's Name and etc.
        \item Pixel Data - actual image contents. each pixel is made up of 2 bytes stored in little endian (reverse byte order).
    \end{itemize}
\end{itemize}

\subsection{DICOM-SEG (.dcm)}
\begin{itemize}
    \item Works like a mask over the whole scan for highlighting selected areas. Is a multi-frame image storing data for each frame and each image.
    \item For sake of storage compression uses BitsAllocated tag to store pixel data as 0 or 1.
\end{itemize}

\subsection{NIfTI (.nii / .nii.gz (compressed))}
\begin{itemize}
    \item stands for Neuroimaging Informatics Technology Initiative. Initially created for brain MRI research, became the standard in medical image analysis.
    \item consists of:
    \begin{itemize}
        \item Fixed header (348 bytes) -  in comparison to DICOM, the NIfTI header is much simpler. It contains no Personal Health Information. Defines matrix dimensions up to 7D, data types and transformation matrices (qform/sform) that map voxel indices into real-world coordinates in millimeters.
        \item Data payload - a continuous binary block containing the entire 3D or 4D volume.
        \item The .nii.gz extension - the exact same file, but compressed with gzip in order to reduce the file size.
    \end{itemize}
\end{itemize}

\subsection{NRRD (.nrrd)}
\begin{itemize}
    \item Nearly Raw Raster Data - lightweight and flexible format designed for storing N-dimensional medical and scientific raster data.
    \item consists of:
    \begin{itemize}
        \item ASCII text header - a human readable text paragraph ending with a blank line storing information such as: type, dimension, axis size, spatial orientation vectors, origin point and data encoding (raw data or compressed using gzip).
        \item Raw data - a multidimensional array stored right after the blank line separating the header section.
    \end{itemize}
    \item purpose: serves as the native working format for 3D slicer and is widely used in 3D visualization. It handles complex data and overlapping segmentation masks really well.
\end{itemize}

\end{document}